\subsection{From transported flag terms to embedded triples}
\label{sec:flag-term-correspondence}

The comparison below concerns the actual double-coset terms, not
a count of free choices of their frame representatives.  It uses the
normal-form arithmetic of the preceding sections.  We separate the
rational normal-form incidence (existence for each relevant term)
from the subsequent integral and weighted comparison.

Write \(F=(e_0,\ldots,e_{r-1},z)\), \(X=\langle e_a\rangle_E\),
\(X_a=\langle e_b:b<a\rangle_E\), and
\(Q_X=X^\perp/X\).  The norm of \(z\) is \(-T\ne0\).
Use \(N_F\) for the isometries acting as one on all successive
isotropic flag quotients and on \(Q_X\).  In particular they
send \(z\) to \(z+X\).  The source row stabilizer \(P_F\)
preserves all \(X_a\) and fixes \(z\) itself.
For a row lattice \(\Lambda\), put
\(N_\Lambda=N_F\cap\operatorname{Aut}(\Lambda)\).
Let \(A_a=\mathfrak d^{-(r-a)}\) be the original leading ideals.
They are conjugation invariant.

\begin{lemma}[Integral movement at a fixed lattice]
\label{lem:integral-frame-movement}
Let \(\Lambda\) be unimodular and let \(F,G\) be two flags with
the same graded vectors, including the last vector modulo \(X\).
Suppose their prescribed leading ideals are \(A_a\), the last
leading ideal is \(\mathcal O_E\), and
\[
 A_a e_a,\ A_a e'_a\subset\Lambda,\qquad z,z'\in\Lambda .
\]
Then there is \(m\in N_\Lambda\) sending every vector of \(F\)
to its corresponding vector of \(G\).
\end{lemma}

\begin{proof}
The successive leading-ideal conditions give
\[
 \Lambda\cap(X+Ez)=\bigoplus_{a<r}A_ae_a\oplus\mathcal O_Ez .
\]
Indeed subtract the last coefficient and then the preceding
coefficients in descending order.  This submodule is saturated,
so its quotient in \(\Lambda\) is torsion free, hence projective
over the Dedekind ring \(\mathcal O_E\).
Each coordinate functional on it extends, by a projective splitting,
to a functional on \(\Lambda\) with values in \(A_a\).
Represent its \(E\)-linear extension as \(h(-,y_a)\).
Thus \(h(e_b,y_a)=\delta_{ba}\), \(h(z,y_a)=0\) and
\(h(\Lambda,y_a)\subset A_a\).  In particular unimodularity
gives \(A_a^{-1}y_a\subset\Lambda\), using conjugation invariance
of \(A_a\).  Integrality also gives
\(h(y_a,y_b)\in A_a A_b\).

Choose the integral trace-one element \(w_0+\bar w_0=1\)
available in the finite-unramified CM extension, and replace
\[
 y_a\longmapsto y_a-w_0\sum_b h(y_a,y_b)e_b .
\]
The pairings of these new vectors vanish: the two cross terms
subtract \(w_0+\bar w_0\) times the old Gram matrix.  All other
pairings just prescribed are unchanged.  The replacement still
has \(A_a^{-1}y_a\subset\Lambda\), by the ideal inclusion for
the Gram entries.  We have split off the unimodular hyperbolic
summand with lattice
\(\bigoplus A_ae_a\oplus\bigoplus A_a^{-1}y_a\);
its rational perpendicular \(C\) contains \(z\).
Repeat for \(G\), obtaining \(y'_a,C'\).  Both complements
identify isometrically with \(Q_X\).  The identification
\(C\to C'\) is orthogonal projection along the new hyperbolic
summand.  It takes the integral lattice of \(C\) into that of
\(C'\), because both hyperbolic summands are unimodular.
Its inverse has the same property.

Match the two hyperbolic bases and use this isometry of
complements.  The resulting isometry preserves \(\Lambda\),
acts identically on \(Q_X\), and acts on \(X\) by the prescribed
unitriangular change.  Since \(z,z'\) are the two lifts of the
same quotient vector in \(C,C'\), it sends \(z\) to \(z'\).
This is the required \(m\).
\end{proof}

\paragraph{Admission of the integral term.}
The split row lattice and the transported frame give the specified
leading ideals and the lattice inclusions in the lemma.  The
source's fixed-congruence condition at every prime dividing \(g\)
says that the localized isotropic coefficient lattice is generated
by the frame itself.  Its admissibility condition gives an integral
upper unitriangular matrix \(\xi\) with
\(G_a\equiv(F\xi)_a\pmod{g\Lambda_{\mathfrak p}}\)
at these primes.  Consequently the following literal condition
holds (the convention is \(X_r=X\), \(Y_r=Y\)):
\[
 x'\in\Lambda,\quad x'-tG_a\in Y_a
 \quad\Longrightarrow\quad
 \exists x\in\Lambda:\ x-tF_a\in X_a,\quad x-x'\in g\Lambda .
 \tag{FC1}
\]
In fact expand \(x'\) in the frame \(G_b\), \(b\le a\).
Its coefficients are in \(A_b\), or in \(\mathcal O_E\)
for the last vector.  Replace \(G_b\) by \((F\xi)_b\).
This gives a vector in \(\Lambda\), since \(A_b\subset A_i\)
for \(i\le b\).  Its leading coefficient is unchanged.
At primes dividing \(g\) the coefficients are integral locally
and each difference is in \(g\Lambda_{\mathfrak p}\).
At all other primes \(g\) is a unit.  Intersection of the
localized lattices proves (FC1) over \(\mathcal O_E\).
Changing the frame by a same-graded change does not change (FC1)
or the leading ideals.

The reverse of (FC1) follows by induction through the flag:
for a given first-flag vector, choose a second-flag vector
with the same leading coefficient (the leading ideals agree);
use (FC1), and subtract the resulting first-flag vector.
The difference belongs to the preceding stage, to which the
induction applies.  Adding the lifts proves the reverse condition.
In particular the primitive submodules
\(\Lambda\cap X\), \(\Lambda\cap Y\) have the same images
modulo \(g\Lambda\); the analogous assertion holds for
\(X+Ez\) and \(Y+Ez'\).  This proves the integral
\texttt{IsAdm} and \texttt{IsTerm} conditions once an adapted
rational normal form is given.

\begin{proposition}[Term-class multiplicity and uniqueness]
\label{prop:flag-term-class-multiplicity}
Fix the source row representative set (one row lattice in each
\(P_F\)-orbit).  A term is a double coset represented by
\((\Lambda,\Lambda',\gamma)\), \(\gamma\Lambda'=\Lambda\),
modulo the two groups \(N_\Lambda,N_{\Lambda'}\).
Let \(G=\gamma F\).  At fixed \(\Lambda\), two frames are
equivalent if \(n\in N_\Lambda\) carries one to the other
up to a same-graded change.  Then every finite set of terms
in one frame class has total reciprocal row weight at most one.
Moreover for terms admitting the rational normal form, the
frame class together with the induced definite-space isometry
determines the term.
\end{proposition}

\begin{proof}
Suppose \(\gamma_1 F,\gamma_2 F\) are equivalent via \(n\).
The isometry \(g_0=\gamma_2^{-1}n\gamma_1\) sends the first
source row lattice to the second, and \(g_0F\) has the
same graded vectors as \(F\).  Both frames at the second
source row satisfy the preceding lemma, since the prescribed
ideal multiples and last vector are transported from their rows.
Choose \(m\in N_{\Lambda'_2}\) with \(mF=g_0F\).
Then \(a=g_0^{-1}m\) fixes \(F\) pointwise and carries
\(\Lambda'_2\) to \(\Lambda'_1\).  The uniqueness of the row
representatives modulo \(P_F\) gives \(\Lambda'_2=\Lambda'_1\).
Write this common row as \(\Lambda'\); the decomposition is
\[
                  \gamma_2=n\gamma_1 a m^{-1}.       \tag{FC2}
\]
Here \(a\in P_F\cap\operatorname{Aut}(\Lambda')\).

For a fixed initial term, assign to each term in the frame class
the induced action of \(a\) on \(Q_X\), choosing one decomposition
(FC2).  This assignment is injective: if the actions for \(a_1,a_2\)
agree, \(a_1^{-1}a_2\) fixes every flag vector and acts as one on
\(Q_X\), hence lies in \(N_{\Lambda'}\).
Equation (FC2) identifies the original double cosets.
The image lies in the actual source group \(\operatorname{autLz}(F,\Lambda')\)
of induced actions, not in a group of lifts to \(V\).
It is finite, because it preserves the quotient lattice, fixes
the nonzero-norm vector \(z\), and is unitary on its totally
definite perpendicular.  A lattice is discrete in the
simultaneous archimedean embedding, so a bounded set of
possible images of generators is finite.  Division by
\(\#\operatorname{autLz}(F,\Lambda')\) proves the weight claim.

For the uniqueness with the definite isometry, write
\(Y_i=\gamma_i X\), \(S_i=(X+Y_i+Ez+E\gamma_i z)^\perp\).
The normal form gives
\[
                 S_i\cap X=S_i\cap Y_i=X\cap Y_i.       \tag{FC3}
\]
Use the fixed definite embedding \(e:W\to V\), where
\(X^\perp\cap z^\perp=eW\oplus X\).
The source's induced isometry \(h_i\) is characterized by
\[
 y_i\in S_i,\quad y_i-e(h_iw)\in X,\quad
 y_i-\gamma_i e(w)\in Y_i\qquad(w\in W).
 \tag{FC4}
\]
Its existence and uniqueness are verified below.
Assume \(h_1=h_2\).  Equation (FC2) and the same-graded
conditions give \(S_2=nS_1\): \(n\) carries the paired flag
spans (including their last vectors) onto one another.
Since \(n^{-1}\) is one on \(Q_X\),
\(y_1-n^{-1}y_2\in X\cap S_1\subset Y_1\).
Also \(m^{-1}e(w)-e(w)\in X\), and \(a\) fixes \(X\).
Thus (FC4) gives
\[
 n^{-1}y_2-\gamma_1 a e(w)\in Y_1,\qquad
                a e(w)-e(w)\in X.
\]
Together with \(az=z\) these show that \(a\) is one on
\(Q_X\); hence \(a\in N_{\Lambda'}\).
Equation (FC2) now gives equality of the terms.
\end{proof}

\paragraph{The two definite quotient identifications.}
For the remaining argument put \(Y=\gamma X\),
\(S=(X+Y+Ez+E\gamma z)^\perp\), \(I=X\cap Y\).
In the admitted normal form the first \(j\) vectors of
each flag have invertible cross pairing, the remaining
isotropic vectors are common up to the scalar \(c\), and
the last vectors are proportional after same-graded changes.
It follows directly from this decomposition that \(S/I\)
is nondegenerate and totally definite.  Each of the two
inclusions into the definite quotient of the corresponding
flag induces an isometry
\[
 S/I\simeq (X+Ez)^\perp/X,\qquad
 S/I\simeq (Y+E\gamma z)^\perp/Y.                    \tag{FC5}
\]
For surjectivity, subtract a linear combination of the
first \(j\) vectors of the corresponding flag, solving the
invertible cross-pairing system.  The tail and last-vector
orthogonalities already hold.  The kernel is (FC3).
The induced map between the two copies of \(W\), using
\(e\) and \(\gamma e\), is exactly \(h=h_\gamma\)
of (FC4); it is an isometry by (FC5).
These statements and their construction are unaffected
by same-graded changes of either flag.

Let
\[
 L_0(\Lambda)=e^{-1}(\Lambda+X),\qquad
 \widetilde L_F=(\Lambda+X)\cap S,\quad
 \widetilde L_G=(\Lambda+Y)\cap S.
\]
Both tilded modules contain the full vector space \(I\);
their quotients by \(I\) are lattices.
Under the first identification in (FC5), they become
\(L_0(\Lambda)\) and \(hL_0(\Lambda')\).
By the earlier unimodular-quotient and primitive
last-vector argument (WF1)--(WF7), both have cyclic
discriminant module \(\mathcal O_E/T\mathcal O_E\).
Define the graph in these quotients by the original lifts:
\[
 (u',u)\in\widetilde L_G\times\widetilde L_F,\quad
 u'-\lambda'\in Y,\quad u-\lambda'-gm\in X,
 \qquad \lambda',m\in\Lambda .                         \tag{FC6}
\]
This is precisely (WF2): absorb the \(X\)-lattice summand
into the second lift.  It contains \(I\times I\).
Denote the corresponding embedded graph in \(W\times W\)
by \(\Gamma_\gamma\).

\begin{lemma}[The embedded graph and its fixed index]
\label{lem:flag-embedded-graph}
Under the preceding conditions \(\Gamma_\gamma\) is the
inverse image of an isometry modulo \(g\) from
\(hL_0(\Lambda')/g hL_0(\Lambda')\) to
\(L_0(\Lambda)/gL_0(\Lambda)\).
The datum
\(d_\gamma=(hL_0(\Lambda'),\Gamma_\gamma)\)
is an embedded datum of Section~\ref{sec:embedded-isometry-genus-count}
with index \(x\) equal to the original \(gI\) of its
normal form, \(1\le x\le\Delta\), and generator bound \(r\).
\end{lemma}

\begin{proof}
For \(u'\in\widetilde L_G\), choose \(\lambda'\in\Lambda\)
with \(u'-\lambda'\in Y\).
For each \(v\in\Lambda\cap(X+Ez)\), admission gives a
congruent vector in \(\Lambda\cap(Y+E\gamma z)\),
which pairs to zero with \(\lambda'\).
Thus \(h(v,\lambda')/g\in\mathcal O_E\).
The submodule \(\Lambda\cap(X+Ez)\) is saturated, so
this functional extends to \(\Lambda\); unimodularity
represents the extension as \(h(-,m)\) for \(m\in\Lambda\).
Then \(\lambda'-gm\) is perpendicular to \(X+Ez\).
By (FC5) subtract an element of \(X\) to obtain
\(u\in S\), proving that the first projection of the graph
is onto.  Reverse admission proves the other projection.
If the first coordinate is zero in \(S/I\), a lift
\(\lambda'\) in (FC6) belongs to \(\Lambda\cap Y\).
Match it modulo \(g\Lambda\) with an element of
\(\Lambda\cap X\).  Then \(u\in g\Lambda+X\), and
\(u/g\in\widetilde L_F\).  This proves the required
kernel in the second coordinate; the first kernel follows
by the reverse argument.  The graph contains both
multiples by \(g\), by taking the corresponding lifts.
Finally the pairing of two graph coordinates equals
the pairing of the chosen lattice lifts: the differences
are in the respective isotropic subspaces, perpendicular
to the graph coordinates.  Their two lattice pairings
are congruent modulo \(g\) by (FC6).  This proves
the asserted isometry of the reduced modules.

In a fixed normal form identify \(S/I\) with its
definite complement \(W_N\) in (US1).  The graph is
then (WF2), with the same diagonal fixed lattice \(K\);
its index is exactly \(x=[L:K]\).  It lies between one
and \(\Delta\) by (WF4).  To check the generator condition
over the ring of integers (not just over \(E\)), map
\(L\) into \(X/((\Lambda\cap X)+I)\) by taking the
\(X\)-part of a lift of \(w\in W_N\cap(\Lambda+X)\).
The map is well-defined and its kernel is precisely
\(M_0=\operatorname{pr}_{W_N}(\Lambda\cap S)\).
Locally the finite image is a submodule of the quotient
of a \(j\)-dimensional fraction-field space by a full
lattice, and so is generated by at most \(j\) elements
over the local DVR (use Smith form inside a bounded
denominator lattice).  A finite torsion module over
\(\mathcal O_E\) with these local bounds is generated
globally by at most \(j\) elements: use its finitely many
primary summands and the Chinese remainder theorem.
As \(M_0\subset K\), the quotient \(L/K\) is generated
by at most \(j\le r\) elements as required.
\end{proof}

\begin{proposition}[Actual row pairs give embedded triples]
\label{prop:flag-pair-embedded-triple}
Suppose two terms \(\gamma,\gamma'\) have the same target
and source row lattices, and satisfy the original
row congruence
\[
 (\gamma'\gamma^{-1}-1)(\Lambda\cap X^\perp)
       \subset g\Lambda+(\Lambda\cap X).
 \tag{FC7}
\]
Then, with the indices of Lemma~\ref{lem:flag-embedded-graph},
\((d_\gamma,d_{\gamma'},h_\gamma h_{\gamma'}^{-1})\)
belongs to the original embedded triple set
\(\mathcal T_r(x,x';L_0(\Lambda),g)\).
\end{proposition}

\begin{proof}
For \(w,v\in W\), (FC4)--(FC6) give the useful
criterion
\[
 (h_\gamma w,v)\in\Gamma_\gamma
 \quad\Longleftrightarrow\quad
 \exists l\in\Lambda':\quad
 e(w)-l\in X,\qquad e(v)-\gamma l\in g\Lambda+X.
 \tag{FC8}
\]
Indeed lift both coordinates to \(S\) using (FC5).
For the first lift use (FC4); its difference from
\(\gamma l\) belongs to \(Y\).  For the second
its difference from \(e(v)\) belongs to \(X\).
This gives both directions of (FC8).

If \(l\) satisfies the first condition in (FC8), then
\(\gamma l\in\Lambda\cap Y^\perp\).
Admission for \(X,Y\) makes the functional
\(h(-,\gamma l)/g\) integral on \(\Lambda\cap X\).
Extend it using saturation and represent it by
\(h(-,m)\), \(m\in\Lambda\), as above.
Thus \(\gamma l-gm\in\Lambda\cap X^\perp\).
Applying (FC7), and using that \(\gamma'\gamma^{-1}\)
preserves \(\Lambda\), proves
\(\gamma'l-\gamma l\in g\Lambda+X\).
The criterion (FC8) is therefore unchanged on replacing
\(\gamma\) by \(\gamma'\), apart from the indicated first
coordinate isometry.  The lattice equality follows from
the common source lattice and (FC5).
These are exactly both isometry conditions (EI2).
\end{proof}

\paragraph{The sum on a prescribed embedded datum.}
For a fixed label \(j,(\delta_l)\) and embedded datum \(d\),
choose at most one term per frame class.
Partition these representatives by the full signed block
matrix and scalar \(c\) of their adapted normal forms.
The bounds (WF12)--(WF13) and the signed block count
give at most \(8^{r+1}f^2\Delta^{2(2+c_E)}\) parts.

Within a part choose a reference normal form \(N_0\).
There are isometries \(\sigma_i\in N_F\) carrying
both frames of \(N_0\) to the corresponding adapted frames
of \(N_i\).  Here is the linear construction.  Match the
first \(j\) vectors of the two flags, the common tail,
the last vector, and the isotropic duals to the tail.
The Gram matrices of these \(2r+1\) independent vectors
are identical, by the full signed block matrix, \(c\)
and \(T\).  Their complements are identified isometrically
with \(eW\) modulo \(X\), by (FC5) and the choice of tail
duals.  Match the complement vectors inducing the same
vectors in that quotient.  This extends to a full isometry
which is one on \(Q_X\), including its last-vector line.
On the first frame it is same-graded, hence belongs to \(N_F\).

Pull back \(\Lambda\) by \(\sigma_i^{-1}\).
The reference first quotient lattice \(\mathcal J\) of (WF1)
is unchanged: \(\sigma_i\) induces the identity on \(Q_X\).
The second definite lattice and graph are also unchanged,
since their embedded data in the fixed \(W\) are precisely
\(d\) under (FC5).  In detail, \(\sigma_i\) transports \(S/I\)
to the corresponding quotient, with the identity on its
image in \(Q_X\); (FC6) transports the graph by that
same identification on both coordinates.
Only \(\mathcal J'\) may vary, and (WF10) allows that
variation with factor \(N_T^2\).

The pulled-back lattices are in distinct orbits of the
common unipotent group for \(N_0\).  Otherwise, conjugating
a common unipotent transporter by the two \(\sigma_i\)
would give an element of \(N_\Lambda\) carrying the
original second frames to one another up to a same-graded
change, contradicting the distinct frame classes.
The weight \(\operatorname{omI}\) is an infimum over adapted
forms, so each weight is at most the weight for its
chosen adapted form, which is unchanged by isometric
transport.  It is constant on a frame class: simultaneous
transport gives a bijection of adapted forms, and preserves
their relative stabilizer indices.  The original normal-form
sum (WF10) now applies within each part.  Multiplying by
the number of parts gives exactly (WF14), i.e.
\(\operatorname{datBd}(k,r,\Delta,D_b,N_T,x)\).

\begin{theorem}[Weighted fixed-row comparison]
\label{thm:weighted-fixed-row-comparison}
Take a finite collection of original relevant ordered term
pairs on a fixed target row, with the two prescribed
labels and indices \(j,j'\).
Suppose the rational adapted normal forms exist for these
terms, with the stated signed block matrix, scalar and
leading-minor label (the original \texttt{NFExC}, \texttt{LabC}
incidence).  Use their admitted row conditions and the
original congruence (FC7).
The weighted sum of the two class weights times the
reciprocal source row weight is at most
\[
 \sum_{1\le x\le\Delta}\sum_{1\le x'\le\Delta'}
 \#\mathcal T_r(x,x';L_0(\Lambda),g)\,
 \operatorname{datBd}(k,r,\Delta,D_b,N_T,x)
 \operatorname{datBd}(k,r,\Delta',D_b',N_T,x').
 \tag{FC9}
\]
The sums range over integers.  There is no extra
automorphism-group or field-degree factor.
\end{theorem}

\begin{proof}
Partition the pairs by the actual triple of
Proposition~\ref{prop:flag-pair-embedded-triple}, including
its isometry \(\eta\).  Inside a triple fibre, fix the
two frame classes.  For each first term there is at most
one second term of the specified class: its isometry
is forced by \(h_{\gamma'}=\eta^{-1}h_\gamma\), and
Proposition~\ref{prop:flag-term-class-multiplicity} gives
uniqueness.  The sum of the reciprocal source row weights
of these first terms is at most one by that proposition.
Both class weights are constant on their classes.
Consequently the fibre sum is at most the product of the
two sums of class weights on the prescribed embedded data.
The preceding \(\operatorname{datBd}\) comparison bounds
these sums.  Sum over triples and indices, using
nonnegativity throughout.
\end{proof}

This proves the lattice, quotient-graph, reciprocal-mass and
fixed-triple fibre comparison.  To invoke it for every term of
the original source, the rational normal-form incidence is supplied by
Theorem~\ref{thm:flag-original-nf-labels}. The remaining
interface is the transport of the original row-label
congruence to (FC7).  The subsequent genus mean,
parameter estimates, zero-Gram sum and generic-return and
finite-column estimates remain distinct steps.
